% !TEX root = ../master.tex
\section{Vom Python-Werkzeug PLY zu Lezer}
\label{subsec:ply_to_lezer}

In den ursprünglichen Python-Implementierungen der Lehrveranstaltung werden zentrale Schritte der Sprachverarbeitung mit der Bibliothek \emph{PLY} (Python Lex-Yacc) umgesetzt. \emph{PLY} ist eine Parsergenerator-Bibliothek für Python, 
die sich an den klassischen Werkzeugen \emph{Lex} und \emph{Yacc} orientiert. 
Diese Werkzeuge werden häufig in Lehrumgebungen eingesetzt, um grundlegende Konzepte der 
Compilertechnik wie Tokenisierung und Parsing praktisch umzusetzen. 
Die Bibliothek wurde von David Beazley entwickelt und ist frei verfügbar 
\autocite{beazley_ply}.

Die grundlegende Architektur eines solchen Sprachprozessors besteht aus zwei zentralen Phasen: der lexikalischen Analyse und der syntaktischen Analyse. Zunächst zerlegt der Lexer den Eingabetext in eine Folge elementarer Token, anschließend analysiert der Parser diese Tokenfolge anhand einer kontextfreien Grammatik und konstruiert daraus eine syntaktische Struktur, etwa einen abstrakten Syntaxbaum (AST) \autocite[S.5ff]{aho2006compilers}.

\emph{PLY} bildet diese beiden Komponenten explizit ab:

\begin{itemize}
    \item \textbf{Lexikalische Analyse (Lexer):}  
    Token werden über reguläre Ausdrücke definiert. Der Lexer erkennt daraus elementare Spracheinheiten wie Zahlen, Operatoren oder Schlüsselwörter.

    \item \textbf{Syntaxanalyse (Parser):}  
    Die Grammatik der Sprache wird als Menge von Produktionsregeln formuliert. \emph{PLY} generiert daraus einen Parser, der auf LR-Parsing-Techniken basiert. Diese Parsingverfahren lesen die Eingabe von links nach rechts und konstruieren eine rechtsableitende Analyse, wodurch auch komplexe Programmiersprachen effizient verarbeitet werden können.
\end{itemize}

Ein charakteristisches Merkmal von \emph{PLY} ist, dass Token-Definitionen und Grammatikregeln direkt im Python-Code definiert werden. Produktionsregeln werden dabei als Funktionen formuliert, deren Docstrings die jeweilige Grammatikregel enthalten. Diese enge Integration erleichtert die prototypische Entwicklung von Parsern im Lehrkontext.

\begin{minted}{python}
tokens = ('NUMBER', 'PLUS')

t_PLUS = r'\+'

def t_NUMBER(t):
    r'\d+'
    t.value = int(t.value)
    return t
\end{minted}

Während \emph{PLY} somit eine klassische Compilerarchitektur innerhalb der Python-Umgebung realisiert, verfolgt das TypeScript-Ökosystem einen anderen Ansatz. Für editornahe Anwendungen und webbasierte Entwicklungsumgebungen wird häufig das Parser-Framework \emph{Lezer} eingesetzt.

\emph{Lezer} verfolgt das Ziel, syntaktische Analysen effizient und inkrementell durchzuführen. Im Gegensatz zu traditionellen Parsergeneratoren wird der Parser so konstruiert, dass er Änderungen am Quelltext lokal verarbeiten kann, ohne den gesamten Syntaxbaum neu berechnen zu müssen \autocite{lezer_docs}.

Die Grammatik wird hierbei nicht mehr als Python-Code formuliert, sondern als deklarative Grammatikbeschreibung in einer eigenen Syntax definiert:

\begin{minted}{typescript}
@top Expression { Add }

Add {
  Expression "+" Expression
}
\end{minted}

Aus dieser Beschreibung generiert \emph{Lezer} einen Parser, der direkt im TypeScript-Ökosystem verwendet werden kann.

Der zentrale Unterschied zwischen beiden Ansätzen liegt somit weniger in der zugrundeliegenden formalen Theorie. Beide basieren auf kontextfreien Grammatiken und klassischen Parsingverfahren, sie Unterscheiden sich vielmehr in der Zielumgebung. Während \emph{PLY} primär für klassische Compilerprototypen und Lehrzwecke in Python konzipiert wurde, ist \emph{Lezer} auf moderne Entwicklungswerkzeuge und inkrementelle Analyse von Programmtexten ausgelegt.

Im Kontext dieser Arbeit dient \emph{PLY} daher als Referenzimplementierung der ursprünglichen Python-Notebooks, während \emph{Lezer} die Grundlage für die Umsetzung der gleichen Sprachverarbeitungsschritte im TypeScript-Ökosystem bildet.

\subsection{Grundlagen von Lezer}
\label{subsec:lezer_basics}

\emph{Lezer} ist ein Parser-Framework, das im Umfeld des CodeMirror-Projekts entwickelt wurde und speziell auf den Einsatz in modernen Entwicklungswerkzeugen ausgelegt ist \autocite{lezer_docs}. Im Gegensatz zu vielen klassischen Parsergeneratoren steht bei \emph{Lezer} nicht die einmalige Analyse eines vollständigen Programms im Vordergrund, sondern die effiziente Verarbeitung sich kontinuierlich ändernder Programmtexte.

Ein zentrales Ziel von \emph{Lezer} besteht darin, die syntaktische Struktur eines Textes auch während der Bearbeitung im Editor schnell verfügbar zu machen. Funktionen wie Syntax-Highlighting, automatische Einrückung oder Fehlermarkierungen benötigen eine konsistente syntaktische Analyse des aktuellen Dokumentzustands. Da sich der Quelltext während der Bearbeitung ständig verändert, muss der Parser in der Lage sein, diese Änderungen effizient zu verarbeiten.

\paragraph{Inkrementelles Parsing}

Eine zentrale Eigenschaft von \emph{Lezer} ist das sogenannte \emph{inkrementelle Parsing}. Dabei wird der Syntaxbaum eines Dokuments nicht bei jeder Änderung vollständig neu berechnet. Stattdessen werden nur jene Teile des Baums aktualisiert, die tatsächlich von der Änderung betroffen sind. 

Dieser Ansatz reduziert den Rechenaufwand erheblich und ermöglicht eine nahezu verzögerungsfreie Aktualisierung der Syntaxstruktur im Editor. Besonders bei größeren Dokumenten oder komplexeren Grammatiken stellt dies einen entscheidenden Vorteil gegenüber klassischen Parserarchitekturen dar \autocite{haverbeke2019lezer}.

\paragraph{Baumrepräsentation der Syntax}

Das Ergebnis der Analyse durch \emph{Lezer} ist eine baumartige Repräsentation der syntaktischen Struktur des Eingabetextes. Diese Struktur enthält sämtliche syntaktischen Elemente, die durch die Grammatik definiert wurden. 

\emph{Lezer} erzeugt dabei zunächst einen sogenannten \emph{Concrete Syntax Tree} (CST). Im Gegensatz zu einem abstrakten Syntaxbaum bleiben hier auch Zwischensymbole und strukturelle Details der Grammatik erhalten. Diese vollständige Darstellung der syntaktischen Struktur ist besonders für editornahe Funktionen relevant, da sie eine präzise Lokalisierung einzelner Sprachelemente ermöglicht.

\paragraph{Integration in das JavaScript- und TypeScript-Ökosystem}

Ein weiterer wichtiger Aspekt von \emph{Lezer} ist seine Integration in das JavaScript- und TypeScript-Ökosystem. Der generierte Parser kann direkt in Webanwendungen oder Entwicklungswerkzeuge (wie z.B in Jupyter-Notebook) eingebunden werden. 

Im Kontext dieser Arbeit wird \emph{Lezer} verwendet, um die syntaktische Analyse der ursprünglichen Python-Lehrbeispiele in eine TypeScript-basierte Implementierung zu übertragen. Die Grammatik der Sprache bildet dabei die formale Grundlage für die Generierung des Parsers, während TypeScript für die Weiterverarbeitung der analysierten Strukturen eingesetzt wird.

Der nächste Abschnitt beschreibt daher die konkrete Definition der Grammatik, die als Grundlage für die spätere Parsergenerierung dient.